\documentclass[10pt,letterpaper]{article}

\usepackage[top=0.3in,bottom=0.3in,left=0.7in,right=0.7in]{geometry}
\usepackage[hidelinks]{hyperref}
\usepackage{enumitem}
\usepackage{titlesec}
\usepackage{xcolor}
\usepackage{ragged2e}

\definecolor{accent}{HTML}{1A3D63}

\pagestyle{empty}
\setlength{\parindent}{0pt}
\RaggedRight

\titleformat{\section}{\large\bfseries\color{accent}}{}{0em}{}[\vspace{-4pt}\titlerule]
\titlespacing{\section}{0pt}{4pt}{2pt}

\setlist[itemize]{leftmargin=14pt, itemsep=0pt, topsep=0pt, parsep=0pt}

% external link: accent-colored so it reads as a link but still line-breaks
\newcommand{\extlink}[2]{\href{#1}{\color{accent}#2}}

\newcommand{\role}[4]{%
  \textbf{#1} \hfill \textit{#3}\\
  \textit{#2} \hfill \textit{#4}\\[0pt]
}

\begin{document}

% ---------- Header ----------
\begin{center}
  {\LARGE \textbf{Cargill Seiveright}}\\[3pt]
  {\large Software Engineer}\\[6pt]
  \small
  Shinjuku, Tokyo, Japan \quad $\vert$ \quad
  \href{mailto:coolgill12@gmail.com}{coolgill12@gmail.com} \quad $\vert$ \quad
  +81 80-3853-9559 \\[0pt]
  \href{https://github.com/gill876}{github.com/gill876} \quad $\vert$ \quad
  \href{https://www.linkedin.com/in/cargill-s-a074b3125/}{linkedin.com/in/cargill-s-a074b3125}
\end{center}

% ---------- Profile ----------
\section*{Profile}
Software engineer in Tokyo working on production Ruby on Rails, with a research background in NLP
and LLM-based data processing (M.Eng., NAIST). At GA technologies I triage production errors for a
real-estate investment service, RENOSY; I brought Claude Code and Snowflake into that work, cutting
investigations from 1--2 days to under an hour, and made the process available to the whole team.
In personal projects I benchmarked vision models to design a receipt-reading budget app that I run in
production and use every week. Aiming to be an engineer who runs AI/LLM products reliably in
production. Works in Japanese with partner teams (JLPT N2).

% ---------- Experience ----------
\section*{Experience}

\role{App Engineer}{GA technologies Co., Ltd.}{2024/11 -- Present}{Roppongi Ichome, Tokyo, Japan}
\textit{RENOSY backend team (2026/05 -- Present, team of 15)}
\begin{itemize}
  \item Triage production and staging errors from Rollbar alerts for the backend behind the
    customer-facing apps (iOS, Android, web) of RENOSY, GA technologies' real-estate investment
    service; the backend integrates with 14 internal and external services across 5 engineering
    teams. Isolate whether the cause is a domain service, infrastructure, or the application, and
    investigate through CloudWatch, CloudTrail, ECS logs, and a Snowflake mirror.
  \item Cut error investigations from 1--2 days to under an hour by being the first on the team to
    connect the Snowflake CLI to Claude Code; packaged the investigation steps as a Claude Code skill
    and published it on the internal marketplace so anyone on the team can use it.
  \item Identified three production data inconsistencies, one of which no one had known about, and
    worked with the domain teams to correct them.
  \item Spotted signs of memory fragmentation in memory monitoring and proposed adopting a tool to isolate
    the cause.
  \item Restructured escalation reports as ``action needed $\rightarrow$ impact $\rightarrow$
    cause'', removing the follow-up questions that used to come back on each escalation.
  \item Wrote and corrected OpenAPI specs for 30+ of the backend's 100+ endpoints, checking each against
    the implementation, as part of a company-wide effort to end duplicated API docs between
    providers and consumers; coordinated with teams including iOS and Android.
  \item Set up a release check that verifies new features on both mobile and web; migrated RSpec
    fixtures to FactoryBot and replaced \texttt{create} with \texttt{build} to shorten CI runs.
\end{itemize}
\textit{Internal learning \& work-management platform (2024/11 -- 2026/04, sole developer)}
\begin{itemize}
  \item Built a full-stack Rails app that tracks employee learning progress and work hours in one
    place, owning basic and detailed design, DB modeling, API/frontend implementation, and tests.
  \item Designed many-to-many models with strict business-rule validations; implemented
    session-based authentication and authorization for nested resources.
  \item Wrote 300+ lines of RSpec\slash FactoryBot tests covering boundary values and edge cases; kept the
    code maintainable with helper extraction and RuboCop; localized the UI with a Japanese locale
    file; automated DB backup\slash restore with shell scripts to cut dev-environment setup effort.
\end{itemize}

\vspace{0pt}
\role{Software Developer}{Artuvic Limited}{2021/03 -- 2022/08}{Kingston, Jamaica (Remote)}
\begin{itemize}
  \item Led frontend and backend enhancements for a nationwide courier platform (AngularJS, Java
    Play Framework), including model design for a new system and CSV-driven API integrations.
  \item Designed and implemented role-based access control and security features; built internal
    admin tools for operations and user management.
  \item Wrote JUnit unit tests for new system components and managed releases via Sourcetree.
\end{itemize}

% ---------- Research \& Publications ----------
\section*{Research \& Publications}
\begin{itemize}\sloppy
  \item \textbf{Master's thesis:} \extlink{https://library.naist.jp/opac/en/book/111663}{Can
    Large Language Models Outperform Rule-based Dataset Converters for Taxonomic Entities?} (NAIST,
    2024). Designed and evaluated an LLM pipeline that converts taxon NER datasets from Standoff to
    CoNLL format, handling the entity-boundary offsets, stray whitespace, and mixed character
    encodings that break rule-based converters; GPT-4o plus post-processing reached F1 $\geq$ 0.99
    on all three datasets.
  \item \textbf{Co-author (2nd author):} \extlink{https://arxiv.org/abs/2403.15430}{Distilling
    Named Entity Recognition Models for Endangered Species from Large Language Models} (arXiv,
    2024). Generated 3,600 synthetic NER and relation-extraction sentences about endangered species
    with GPT-4 to distill into small BERT models; responsible for prompt design for data generation,
    fact-checking against the IUCN Red List, and writing part of the paper.
\end{itemize}

% ---------- Projects ----------
\section*{Projects}
\begin{itemize}\sloppy
  \item \extlink{https://kechi-plus-g6bg.web.app}{\textbf{Kechi+}} (2026/10 -- Present,
    used weekly): a budget app that reads Japanese receipts from a photo (items, prices, tax
    category) and compares unit prices by store and month. Planned it and made the design decisions;
    built and operated it with Claude Code.
    \begin{itemize}
      \item Benchmarked 11 vision models on GMI Cloud against my own receipts for accuracy, cost,
        and speed; chose a fast model for first reads, with a stronger model re-reading only when
        the totals don't add up.
      \item Trunk-based development; GitHub Actions deploys to Cloud Run after lint and tests pass,
        then smoke-tests the live site; keyless auth via Workload Identity Federation. Monitoring
        with Sentry (errors, traces, per-release tracking), Grafana dashboards, SLOs, and a runbook.
      \item Python (FastAPI), Next.js/TypeScript, PostgreSQL (Cloud SQL), Cloud Run, Firebase
        Authentication.
    \end{itemize}
  \item \extlink{https://github.com/gill876/trunk-development-feature-flag}{\textbf{Flagship}}
    (2026): a lab for trunk-based development, feature flags, and canary releases. Planned it and
    made the design decisions; built and operated it with Claude Code.
    \begin{itemize}
      \item Every merge to main deploys a new Cloud Run revision at 0\% traffic, released through a
        0\% $\rightarrow$ 10\% $\rightarrow$ 100\% canary.
      \item Feature flags built on Rails configuration without a gem: segment targeting (dev, alpha,
        beta), sticky per-user percentage rollout, and a kill switch.
      \item Because old and new revisions serve side by side during a canary, CI uses oasdiff to
        block breaking OpenAPI changes.
      \item 12 bots browse as users, and a dashboard shows which revision and flag state served
        each one; monitored with Sentry, Grafana, and SLOs. Rails 8.1, Cloud Run, Firestore, GitHub
        Actions (Brakeman, bundler-audit, RuboCop).
    \end{itemize}
  \item \extlink{https://zenn.dev/hackathons/google-cloud-japan-ai-hackathon-vol5/projects/bf01c70153667d}{\textbf{5th
    Agentic AI Hackathon with Google Cloud}} (2026): planned and built AsobiAgent solo, an AI agent
    that plans group outings in LINE groups. Owned the planning and design decisions and directed
    Claude Code and Antigravity to implement it. Chose not to let the AI book restaurants by phone,
    given the corporate contract it requires and the burden on restaurants; a person places the call
    with one tap instead. Vertex AI (Gemini), Cloud Run, Firestore, Secret Manager, LINE (LIFF),
    Ekispert API, Hot Pepper API.
  \item Competed in the ACM ICPC as an undergraduate.
\end{itemize}

% ---------- Internships ----------
\section*{Internships}

\role{Research Intern, NLP}{Fujitsu Ltd.}{2023/09 (3 weeks)}{Kanagawa, Japan (Remote)}
\begin{itemize}
  \item Surveyed tokenization and embedding methods for low-resource NLP domains, focused on
    endangered-species corpora, and presented technical recommendations to Fujitsu's NLP Research
    team.
\end{itemize}

\vspace{0pt}
\role{Software Intern}{The Palisadoes Foundation -- Calico Challenge}{2020/06 -- 2020/08}{Kingston, Jamaica (Remote)}
\begin{itemize}
  \item Implemented encrypted communication between a central API and distributed data-collection
    agents for an open-source monitoring system (Pattoo) in Python, and built a React/Tailwind CSS
    admin portal for managing users, agents, and data pipelines.
\end{itemize}

% ---------- Education ----------
\section*{Education}
\textbf{Nara Institute of Science and Technology (NAIST)} \hfill 2022/10 -- 2024/09\\
\textit{Master of Engineering, Natural Language Processing} \hfill \textit{Nara, Japan}\\
MEXT (Monbukagakusho) Scholarship recipient, Japanese Embassy recommendation\\[0pt]
\textbf{The University of the West Indies, Mona} \hfill 2018/08 -- 2021/09\\
\textit{Bachelor of Science, Computer Science (Minor: Japanese)} \hfill \textit{Kingston, Jamaica}

% ---------- Skills ----------
\section*{Skills}
\textbf{Languages:} Ruby, Python, TypeScript, Java, JavaScript, SQL \\
\textbf{Frameworks:} Ruby on Rails, FastAPI, Next.js, React, AngularJS, Flask \\
\textbf{Cloud \& Data:} AWS (ECS, CloudWatch, CloudTrail), Google Cloud (Cloud Run, Cloud SQL,
Firestore), Docker, PostgreSQL, Snowflake \\
\textbf{Observability \& CI:} Sentry, Grafana, Rollbar, GitHub Actions \\
\textbf{Practices/Tools:} RSpec, FactoryBot, TDD, OpenAPI, Claude Code, trunk-based development,
feature flags, canary releases, Git \\
\textbf{Spoken Languages:} English (native), Japanese (JLPT N2, business working proficiency)

\end{document}
